\chapter{Магнитные силы, векторные потенциалы и принцип действия}

\lectureoverview{Электростатические силы укладываются в привычную консервативную схему, но магнитные силы зависят от скорости. Чтобы включить их в лагранжеву механику, мы сначала развиваем необходимый векторный анализ, используем тот факт, что магнитное поле не имеет источников, чтобы ввести векторный потенциал, рассматриваем его калибровочную свободу, а затем строим магнитный член в действии.}

\section{Новый вид силы}

Предшествующее обсуждение момента импульса, гироскопов и скобок Пуассона свело воедино несколько тем. В частности, вращающееся заряженное тело в магнитном поле прецессирует почти так же, как гироскоп. Теперь мы переходим от этой аналогии к движению отдельной заряженной частицы.

Электростатические силы не требуют новой механики. Если \(U(\mathbf x)\) — потенциальная энергия, то
\begin{equation}
  \mathbf F=-\nabla U(\mathbf x),
\end{equation}
и обычный довод дает закон сохранения суммы кинетической и потенциальной энергий. Для частицы с зарядом \(e\) в электростатическом потенциале \(V(\mathbf x)\),
\begin{equation}
  U(\mathbf x)=eV(\mathbf x),
  \qquad
  \mathbf F_{\mathrm{el}}=-e\nabla V.
\end{equation}
Магнитные силы имеют другую форму. Они зависят от скорости частицы, поэтому их нельзя ввести просто добавлением еще одной обычной потенциальной энергии. Тем не менее они следуют из принципа действия и допускают как лагранжево, так и гамильтоново описание. Чтобы прийти к этому описанию, нужен короткий, но существенный экскурс в векторный анализ.

\section{Поля и локальные операции}

\subsection{Что фиксирует поле}

Поле сопоставляет физическую величину каждой точке пространства и, вообще говоря, каждому моменту времени. Температура и давление — скалярные поля. Скорость ветра — векторное поле. Электрическое и магнитное поля — векторные поля. В других теориях скалярные поля встречаются даже на микроскопических масштабах, тогда как гравитация описывается тензорным полем. В настоящем обсуждении мы опускаем всякую зависимость от времени и записываем поля как функции \(\mathbf x\).

Операциональный смысл поля состоит в локальности. Представьте, что некоторая область скрыта внутри мешка. Чтобы измерить электрическое поле в этой области, поместите туда малый пробный заряд и наблюдайте его ускорение. Сила определяется полем в той точке, где находится заряд, а не отдельным обращением к каждому далекому источнику. То же относится и к магнитному полю и движущемуся пробному заряду. Позднее векторный потенциал заставит нас уточнить эту простую картину, но начинать следует именно с нее.

\subsection{Оператор дифференцирования}

Производные превращают поля в другие поля. Мы объединяем три пространственные производные в
\begin{equation}
  \nabla
  =
  \left(
    \frac{\partial}{\partial x},
    \frac{\partial}{\partial y},
    \frac{\partial}{\partial z}
  \right).
\end{equation}
Символ читается как \emph{дель} или \emph{набла}. Его вид перевернутого треугольника легко запоминается, но его подлинное значение состоит в том, что формально он ведет себя как вектор, компонентами которого являются дифференциальные операторы.

Это обычное трехмерное физическое пространство, не фазовое пространство и не конфигурационное пространство многих частиц. Более того, \(\nabla\) — не векторное поле, которое может существовать само по себе. Оно должно действовать на что-то. Если \(T(\mathbf x)\) — скалярное поле температуры, то
\begin{equation}
  \nabla T
  =
  \left(
    \frac{\partial T}{\partial x},
    \frac{\partial T}{\partial y},
    \frac{\partial T}{\partial z}
  \right)
\end{equation}
есть его градиент, подлинное векторное поле.

\begin{figure}[htbp]
\centering
\includegraphics[width=0.82\textwidth]{lecture_09_figure_01.png}
\caption{На доске \(\nabla\) вводится как вектор операторов пространственных производных, после чего его начинают применять к скалярному полю \(T(\mathbf x)\).}
\lectureframe{00:11:30}
\end{figure}

Скаляр — это не просто величина, записанная одной компонентой. Это величина, значение которой не зависит от ориентации координатных осей. Стандартный пример — температура. Для данного векторного поля \(\mathbf E\) локальное выражение \(\mathbf E\cdot\mathbf E\) тоже является скаляром. Напротив, построение потенциала интегрированием поля — операция нелокальная; она использует значения вдоль пути.

\subsection{Дивергенция}

Для векторного поля \(\mathbf A=(A_x,A_y,A_z)\) скаляр
\begin{equation}
  \nabla\cdot\mathbf A
  =
  \frac{\partial A_x}{\partial x}
  +\frac{\partial A_y}{\partial y}
  +\frac{\partial A_z}{\partial z}
\end{equation}
является его дивергенцией. Говоря на языке гидродинамики, положительная дивергенция измеряет локальную тенденцию течь наружу, тогда как отрицательная измеряет тенденцию сходиться внутрь. Источник жидкости создает положительную дивергенцию; сток создает отрицательную. Эта картина качественна, но она сделает наглядным отсутствие источников у магнитного поля.

\section{Векторные произведения, символы эпсилон и ротор}

\subsection{Векторное произведение в компонентах}

Для двух векторов \(\mathbf V\) и \(\mathbf A\) одна из компонент векторного произведения имеет вид
\begin{equation}
  (\mathbf V\times\mathbf A)_x
  =
  V_yA_z-V_zA_y.
\end{equation}
Остальные две получаются циклической заменой \(x\to y\to z\to x\). Полезно заменить \(x,y,z\) индексами \(1,2,3\). Символ Кронекера равен
\begin{equation}
  \delta_{ij}
  =
  \begin{cases}
    1,&i=j,\\
    0,&i\ne j,
  \end{cases}
\end{equation}
а символ Леви--Чивиты определяется как
\begin{equation}
  \epsilon_{123}=\epsilon_{231}=\epsilon_{312}=1,
  \qquad
  \epsilon_{132}=\epsilon_{321}=\epsilon_{213}=-1,
\end{equation}
причем \(\epsilon_{ijk}=0\), если совпадают любые два индекса. Циклический обход по циферблату с тремя делениями дает знаки плюс; обход в обратном направлении дает знаки минус. Перестановка любой пары индексов меняет знак, так что \(\epsilon_{ijk}\) полностью антисимметричен.

Мы также используем соглашение о суммировании: индекс, повторяющийся в одном члене, суммируется от \(1\) до \(3\). Тогда три уравнения для векторного произведения принимают вид
\begin{equation}
  (\mathbf V\times\mathbf A)_i
  =
  \epsilon_{ijk}V_jA_k.
\end{equation}
Например, при \(i=1\) остаются только \((j,k)=(2,3)\) и \((3,2)\):
\begin{equation}
  (\mathbf V\times\mathbf A)_1
  =
  \epsilon_{123}V_2A_3
  +\epsilon_{132}V_3A_2
  =
  V_2A_3-V_3A_2.
\end{equation}
Следовательно, компактная запись в точности содержит привычное покомпонентное правило.

\subsection{Вектор под видом антисимметричной матрицы}

Существует и второе представление:
\begin{equation}
  V_iA_j-V_jA_i
  =
  \epsilon_{ijk}(\mathbf V\times\mathbf A)_k.
\end{equation}
Левая часть имеет два индекса и антисимметрична при \(i\leftrightarrow j\). В виде матрицы она имеет вид
\begin{equation}
  M_{ij}
  =
  \begin{pmatrix}
    0&M_{12}&M_{13}\\
    -M_{12}&0&M_{23}\\
    -M_{13}&-M_{23}&0
  \end{pmatrix}.
\end{equation}
Диагональные элементы равны нулю, а нижний треугольник определяется верхним. Поэтому имеется три независимые компоненты.

\begin{figure}[htbp]
\centering
\includegraphics[width=0.82\textwidth]{lecture_09_figure_02.png}
\caption{Одна компонента векторного произведения и схематический антисимметричный массив \(3\times3\), используемый для подсчета его независимых компонент. Пометки в клетках служат мнемоникой, а не дополнительными символами.}
\lectureframe{00:34:20}
\end{figure}

\begin{classroomqa}
\audiencequestion{Имеет ли векторное произведение, или отождествление с другим вектором, смысл вне трех измерений?}
\lecturerresponse{Антисимметричный объект с двумя индексами имеет смысл при любом числе измерений, но его отождествление с вектором является особенностью трех измерений. Вектор в трех измерениях имеет три компоненты, и антисимметричная матрица \(3\times3\) имеет три независимые компоненты. В четырех измерениях соответствующие числа равны 4 и 6, поскольку \(4(4-1)/2=6\). Антисимметричный объект по-прежнему существует и естественным образом рассматривается как 2-форма, но он не является еще одним четырехвектором.}
\end{classroomqa}

\subsection{Ротор}

Замена \(\mathbf V\) на \(\nabla\) определяет ротор:
\begin{equation}
  (\nabla\times\mathbf A)_i
  =
  \epsilon_{ijk}\frac{\partial A_k}{\partial x_j}.
\end{equation}
В обычных компонентах,
\begin{align}
  (\nabla\times\mathbf A)_x
  &=
  \frac{\partial A_z}{\partial y}
  -\frac{\partial A_y}{\partial z},\\
  (\nabla\times\mathbf A)_y
  &=
  \frac{\partial A_x}{\partial z}
  -\frac{\partial A_z}{\partial x},\\
  (\nabla\times\mathbf A)_z
  &=
  \frac{\partial A_y}{\partial x}
  -\frac{\partial A_x}{\partial y}.
\end{align}
Тогда как дивергенция измеряет локальное растекание, ротор измеряет локальную циркуляцию. У жидкости, вращающейся вокруг оси, ротор вдоль этой оси отличен от нуля. Эта картина особенно наглядна для полей скорости, хотя определение применимо к любому векторному полю.

\section{Два структурных тождества}

Почти вся математическая работа сводится к двум фактам:
\begin{equation}
  \nabla\cdot(\nabla\times\mathbf A)=0,
  \qquad
  \nabla\times(\nabla s)=0.
\end{equation}
Первое следует из того, что вторые производные симметричны, тогда как \(\epsilon_{ijk}\) антисимметричен:
\begin{equation}
  \nabla\cdot(\nabla\times\mathbf A)
  =
  \epsilon_{ijk}\partial_i\partial_j A_k
  =0.
\end{equation}
Второе следует из равенства смешанных частных производных:
\begin{equation}
  \bigl[\nabla\times(\nabla s)\bigr]_i
  =
  \epsilon_{ijk}\partial_j\partial_k s
  =0.
\end{equation}

\begin{figure}[htbp]
\centering
\includegraphics[width=0.80\textwidth]{lecture_09_figure_03.png}
\caption{Первое тождество на доске: дивергенция любого ротора равна нулю.}
\lectureframe{00:38:26}
\end{figure}

Это легкие направления, подходящие для обещанного на лекции домашнего задания или внезапной контрольной. Обратные утверждения требуют большей аккуратности. Локально или, глобально, на подходящей односвязной или стягиваемой области,
\begin{align}
  \nabla\cdot\mathbf F=0
  &\quad\Longrightarrow\quad
  \mathbf F=\nabla\times\mathbf A,\\
  \nabla\times\mathbf F=0
  &\quad\Longrightarrow\quad
  \mathbf F=\nabla s.
\end{align}
Без оговорки об области дыры и другие глобальные топологические особенности могут препятствовать этим выводам. В лекции используется локальная форма, и именно она нужна ниже.

\section{Почему магнитное поле является ротором}

\subsection{Отсутствие наблюдаемых магнитных источников}

Магнитное поле \(\mathbf B\) определяется операционально через его действие на заряженные частицы и токи. В обычном электромагнетизме без магнитных монополей,
\begin{equation}
  \nabla\cdot\mathbf B=0.
\end{equation}
Это отличается от электрического поля, дивергенция которого пропорциональна плотности электрического заряда. Электрические заряды являются источниками и стоками электрического потока. Никакого изолированного источника или стока магнитного потока не наблюдалось.

Магнитные монополи изменили бы это уравнение. Теории единого поля нередко дают теоретические основания ожидать существования таких объектов, но, если они существуют, их масса может быть слишком велика, чтобы их можно было породить в доступных экспериментах. Поэтому в настоящей теории плотность магнитного заряда принимается равной нулю.

\begin{classroomqa}
\audiencequestion{Что такое магнитный монополь? У обычных магнитов уже есть северный и южный полюса.}
\lecturerresponse{Изолированный монополь был бы подлинным источником или стоком магнитного поля, аналогично электрическому заряду. Обычный стержневой магнит устроен иначе. Линии поля, выходящие вблизи его северного полюса, изгибаются в пространстве и возвращаются через южный полюс, продолжаясь внутри материала магнита. То, что снаружи выглядит как источник и сток, на самом деле соединено, как жидкость, текущая по скрытой трубе. Никакого суммарного рождения или уничтожения магнитного потока нет.}
\end{classroomqa}

\begin{figure}[htbp]
\centering
\includegraphics[width=0.76\textwidth]{lecture_09_figure_04.png}
\caption{На доске противопоставлены кажущиеся северный и южный источники, при этом их силовые линии соединены через магнит, что подчеркивает: магнитный поток там не начинается и не заканчивается.}
\lectureframe{00:44:56}
\end{figure}

Поскольку \(\mathbf B\) бездивергентно, обратная теорема позволяет записать, по крайней мере локально,
\begin{equation}
  \boxed{\mathbf B=\nabla\times\mathbf A}.
\end{equation}
Поле \(\mathbf A\) — это векторный потенциал. Поначалу это может выглядеть как ненужное переописание измеримого поля \(\mathbf B\). Его назначение станет ясным лишь тогда, когда мы построим действие.

\subsection{Векторный потенциал не единственен}

Если \(\mathbf A\) — один из векторных потенциалов, то
\begin{equation}
  \mathbf A'=\mathbf A+\nabla s
\end{equation}
— тоже векторный потенциал, поскольку
\begin{equation}
  \nabla\times\mathbf A'
  =
  \nabla\times\mathbf A+\nabla\times(\nabla s)
  =
  \mathbf B.
\end{equation}
Такая замена называется \emph{калибровочным преобразованием}. Магнитное поле калибровочно-инвариантно.

\begin{figure}[htbp]
\centering
\includegraphics[width=0.80\textwidth]{lecture_09_figure_05.png}
\caption{На доске рядом записаны определяющее соотношение \(\mathbf B=\nabla\times\mathbf A\) и калибровочный сдвиг \(\mathbf A'=\mathbf A+\nabla s\).}
\lectureframe{01:03:01}
\end{figure}

Его произносят \emph{gauge}, а не \emph{gouache} и не \emph{gouge}--различие, на котором в лекции с притворной строгостью настаивают. Само слово историческое. Оно возникло в ранней попытке связать эту свободу со сравнением эталонов длины в разных местах; исходная интерпретация не сохранилась, а название осталось.

\begin{classroomqa}
\audiencequestion{В этом обсуждении калибровочное преобразование означает именно \(\mathbf A\mapsto\mathbf A+\nabla s\)?}
\lecturerresponse{Да. Именно это и есть полное электромагнитное преобразование, используемое здесь. Калибровочные преобразования в теории Янга--Миллса и других теориях поля обобщают ту же идею, но в электромагнетизме требуется только добавление градиента скалярной функции.}
\end{classroomqa}

Неединственность также проясняет, что именно непосредственно измеримо. Заряженная частица в малой области реагирует на локальное \(\mathbf B\), и все векторные потенциалы, связанные калибровочным преобразованием, предсказывают там одну и ту же силу. Не существует локального эксперимента такого классического типа, который приписывал бы \(\mathbf A\) единственное значение.

\begin{classroomqa}
\audiencequestion{Является ли добавление градиента полным описанием этой неединственности?}
\lecturerresponse{Локально, в односвязной области, да. Если \(\nabla\times\mathbf A'=\nabla\times\mathbf A\), то \(\nabla\times(\mathbf A'-\mathbf A)=0\), так что разность является градиентом. Глобально многосвязная область может вносить дополнительные тонкости. Добавление постоянного вектора уже включено в это описание, поскольку постоянный вектор является градиентом линейной скалярной функции. Добавление константы к самой \(s\) ничего не меняет, потому что её градиент исчезает.}
\end{classroomqa}

\begin{classroomqa}
\audiencequestion{Если у калибровочной свободы есть лежащая в основе физическая причина, то в чём она состоит?}
\lecturerresponse{Лекция откладывает эту более глубокую интерпретацию. Полезный первый шаг — работать с этой свободой и увидеть, как она входит в действие. Когда формализм станет привычным, будет легче распознать, почему калибровочные описания незаменимы.}
\end{classroomqa}

\section{Одно магнитное поле, много потенциалов}

Рассмотрим простейший возможный пример,
\begin{equation}
  \mathbf B=b\,\hat{\mathbf z},
  \qquad
  B_x=B_y=0,
  \qquad
  B_z=b,
\end{equation}
где \(b\) — константа. Один из векторных потенциалов имеет вид
\begin{equation}
  \mathbf A^{(1)}=(0,bx,0).
\end{equation}
В самом деле,
\begin{equation}
  (\nabla\times\mathbf A^{(1)})_z
  =
  \frac{\partial A_y^{(1)}}{\partial x}
  -\frac{\partial A_x^{(1)}}{\partial y}
  =b,
\end{equation}
а остальные компоненты ротора равны нулю.

Этот выбор, по-видимому, выделяет направление \(y\), хотя физическое поле обладает вращательной симметрией в плоскости \(x\)-\(y\). Второй выбор делает произвольность явной:
\begin{equation}
  \mathbf A^{(2)}=(-by,0,0).
\end{equation}
Он даёт
\begin{equation}
  (\nabla\times\mathbf A^{(2)})_z
  =
  0-\frac{\partial(-by)}{\partial y}
  =b.
\end{equation}
Два потенциала отличаются на
\begin{equation}
  \mathbf A^{(2)}-\mathbf A^{(1)}
  =
  (-by,-bx,0)
  =
  \nabla(-bxy).
\end{equation}
Следовательно, они калибровочно эквивалентны. Вращательно-симметричный выбор,
\begin{equation}
  \mathbf A^{(\mathrm{sym})}
  =
  \frac12\,\mathbf B\times\mathbf r
  =
  \left(-\frac{by}{2},\frac{bx}{2},0\right),
\end{equation}
— ещё один представитель того же семейства.

\begin{figure}[htbp]
\centering
\includegraphics[width=0.80\textwidth]{lecture_09_figure_06.png}
\caption{На двух фрагментах доски записаны варианты \(A_y=bx\) и \(A_x=-by\), оба дающие одно и то же однородное поле \(B_z=b\).}
\lectureframe{01:07:20}
\end{figure}

\begin{classroomqa}
\audiencequestion{Является ли это той же неоднозначностью, что и при различной расстановке катушек с током для создания поля?}
\lecturerresponse{Нет. Разные распределения тока — это разные физические источники, и они могут создавать в других местах разные поля, даже если в одной области совпадают. Калибровочная свобода остаётся и после того, как магнитное поле задано всюду: одно и то же физическое \(\mathbf B\) всё ещё имеет много векторных потенциалов.}
\end{classroomqa}

\begin{classroomqa}
\audiencequestion{Можно ли представлять добавление градиента как растяжение или сжатие без добавления закрутки?}
\lecturerresponse{Это полезная геометрическая интуиция. Ротор выявляет локальную закручивающую часть, а чистый градиент имеет нулевой ротор. Это утверждение в точности выражается равенством \(\nabla\times\nabla s=0\).}
\end{classroomqa}

\begin{classroomqa}
\audiencequestion{Существует ли фиксированное соотношение между направлением \(\mathbf A\) и направлением \(\mathbf B\)?}
\lecturerresponse{Никакого калибровочно-независимого соотношения не существует. В приведённых выше примерах просто так получилось, что \(A_z=0\), но к \(A_z\) можно было бы прибавить константу, не изменив ни одной производной. Тогда \(\mathbf A\) уже не был бы перпендикулярен \(\mathbf B\). В этом классическом обсуждении инвариантное физическое содержание имеет только ротор.}
\end{classroomqa}

В электростатике есть более слабый аналог. Поскольку \(\mathbf E=-\nabla V\), замена \(V\) на \(V+C\) не изменяет \(\mathbf E\). Термин \emph{векторный потенциал} был выбран отчасти по аналогии с этой свободой в обычном скалярном потенциале: потенциал можно изменить так, что физическое поле останется тем же.

\section{Сила Лоренца}

Теперь мы можем сформулировать закон силы, который должен воспроизвести принцип действия. В используемых здесь соглашениях
\begin{equation}
  \boxed{
  \mathbf F
  =
  e\left(
    \mathbf E+\frac{\mathbf v}{c}\times\mathbf B
  \right)
  =
  -e\nabla V+\frac{e}{c}\mathbf v\times\mathbf B
  }.
\end{equation}
Символ \(e\) — это электрический заряд частицы, \(\mathbf v=\dot{\mathbf x}\) — её скорость, а \(V\) — скалярный потенциал. Эта запись не должна побуждать нас спутать \(V\) с \(\mathbf v\). В системе SI \(V\) измеряется в вольтах, а \(e\) — в кулонах; показанный множитель \(1/c\) относится к соглашению о единицах, используемому в лекции. Существенный структурный момент не зависит от выбора единиц.

\begin{figure}[htbp]
\centering
\includegraphics[width=0.82\textwidth]{lecture_09_figure_07.png}
\caption{Электрический и магнитный члены силы Лоренца показаны вместе. Магнитный вклад пропорционален \(\mathbf v\times\mathbf B\).}
\lectureframe{00:58:54}
\end{figure}

Электрический член консервативен и уже знаком. Магнитный член здесь новый. Он обращается в нуль для покоящейся частицы и перпендикулярен и \(\mathbf v\), и \(\mathbf B\). Следовательно,
\begin{equation}
  \mathbf v\cdot\mathbf F_{\mathrm{mag}}
  =
  \frac{e}{c}\mathbf v\cdot(\mathbf v\times\mathbf B)
  =0.
\end{equation}
Статическое магнитное поле искривляет траекторию, но не совершает над частицей работы. Это не похоже на вязкое сопротивление, которое направлено противоположно скорости и отнимает кинетическую энергию.

Множитель \(v/c\) также показывает, почему при сопоставимых масштабах полей в этих соглашениях магнитные эффекты для медленно движущихся частиц часто меньше электрических. Электроны достаточно легки, чтобы достигать больших скоростей, поэтому магнитное отклонение по-прежнему легко наблюдать.

Исторически силы, действующие на проводники с током, были известны раньше, чем была выделена сила, действующая на отдельный заряд. Лоренц понял, что силу, действующую на проводник, можно объяснить, приписав скоростезависимый член каждому движущемуся заряду. Полное выражение, включающее электрическую и магнитную части, теперь называется законом силы Лоренца, хотя в нынешнем обсуждении новым является именно магнитный член.

\begin{classroomqa}
\audiencequestion{Есть ли у всего выражения отдельное обозначение, или силой Лоренца называется только магнитная часть?}
\lecturerresponse{Всё выражение — это просто сила \(\mathbf F=m\mathbf a\), и обычно именно его называют силой Лоренца. Магнитный член — часть, исторически наиболее непосредственно связанная с Лоренцем. Электрические и магнитные поля глубоко связаны относительностью, но это объединение не нужно для данной механической задачи.}
\end{classroomqa}

\subsection{Аналогия с силой Кориолиса}

Мы уже встречали векторное произведение скорости с фиксированным вектором. В системе отсчёта, вращающейся с угловой скоростью \(\boldsymbol\omega\), силу Кориолиса можно записать в виде
\begin{equation}
  \mathbf F_{\mathrm{Cor}}
  =
  2m\,\mathbf v\times\boldsymbol\omega.
\end{equation}
Она имеет ту же алгебраическую форму, что и магнитная сила:
\begin{equation}
  2m\boldsymbol\omega
  \quad\longleftrightarrow\quad
  \frac{e}{c}\mathbf B.
\end{equation}
Поэтому однородное магнитное поле, перпендикулярное плоскости, искривляет траектории заряженных частиц почти так же, как вращающаяся система отсчёта отклоняет движущиеся тела.

\begin{figure}[htbp]
\centering
\includegraphics[width=0.82\textwidth]{lecture_09_figure_08.png}
\caption{Слева остаётся сила Лоренца, а на доске добавляется соответствующая форма Кориолиса \(2m\,\mathbf v\times\boldsymbol\omega\).}
\lectureframe{01:02:08}
\end{figure}

\begin{classroomqa}
\audiencequestion{Силу Кориолиса называют фиктивной. Не является ли она менее реальной, чем магнитная сила, потому что зависит от скорости?}
\lecturerresponse{Зависимость от скорости здесь несущественна. Человек на вращающейся платформе может удерживать движущийся предмет и ощущать силу Кориолиса. Ее называют фиктивной, потому что в инерциальной системе отсчета она исчезает, тогда как вращающаяся система отсчета определяется именно наличием таких инерциальных членов. Слово \emph{реальная} часто менее полезно, чем указание, какая система отсчета и какое уравнение используются.}
\end{classroomqa}

Эту аналогию можно развить в небольшой мысленный эксперимент. На вращающейся платформе подходящее электростатическое поле может компенсировать центробежный член для выбранной частицы, тогда как подходящее однородное магнитное поле может компенсировать член Кориолиса. Такая компенсация работает только для выбранного отношения заряда к массе: электрические и магнитные силы масштабируются с \(e\), тогда как инерциальные силы масштабируются с \(m\). Конфигурация поля, настроенная на электроны, не компенсировала бы те же члены для частиц с другим \(e/m\).

\section{Построение действия для заряженной частицы}

\subsection{Стационарность действия при фиксированных концах}

Теперь мы задаем центральный вопрос: можно ли вывести силу Лоренца из принципа действия? Ответ положителен, но векторный потенциал незаменим. Лагранжиан, локально построенный только из частицы и \(\mathbf B\), не дает требуемой магнитной силы. Вместо того чтобы доказывать это отрицательное утверждение, мы сразу построим действие, которое работает.

Для траектории \(\mathbf x(t)\) запишем
\begin{equation}
  S[\mathbf x]
  =
  \int_{t_i}^{t_f}L(\mathbf x,\dot{\mathbf x},t)\,dt.
\end{equation}
В лекции для действия и векторного потенциала используется одна и та же рукописная буква, но мы будем обозначать действие через \(S\), а векторный потенциал через \(\mathbf A\). Физическая траектория делает \(S\) стационарным относительно вариаций, удерживающих оба конца фиксированными. Традиционное выражение \emph{наименьшее действие} удобно, хотя термин «стационарное действие» точнее, поскольку экстремум не обязан быть минимумом.

Фиксированность концов существенна. Можно представить, что концы пути прибиты к доске, а внутренняя часть слегка деформируется. Это траектория в обычной механике, а не объект теории струн; образ струны служит лишь способом наглядно представить вариацию.

\subsection{Что может входить в лагранжиан?}

При отсутствии всех полей лагранжиан равен кинетической энергии,
\begin{equation}
  L_0=\frac{m}{2}\dot{\mathbf x}^{\,2}.
\end{equation}
Электростатический член тоже стандартен:
\begin{equation}
  L_{\mathrm{el}}=-eV(\mathbf x).
\end{equation}
Магнитный член должен удовлетворять нескольким условиям:
\begin{enumerate}
  \item быть пропорциональным заряду \(e\);
  \item содержать множитель \(1/c\) в единицах, используемых в лекции;
  \item содержать скорость;
  \item быть скаляром относительно вращений.
\end{enumerate}

Соблазнительный пробный член \(\mathbf B\cdot\dot{\mathbf x}\) является скаляром и содержит скорость, но приводит к неверному уравнению движения. Правильный член использует векторный потенциал:
\begin{equation}
  L_{\mathrm{mag}}
  =
  \frac{e}{c}\mathbf A(\mathbf x)\cdot\dot{\mathbf x}
  =
  \frac{e}{c}A_i(\mathbf x)\dot x_i.
\end{equation}
Таким образом, полный не зависящий от времени электромагнитный лагранжиан имеет вид
\begin{equation}
  \boxed{
  L
  =
  \frac{m}{2}\dot{\mathbf x}^{\,2}
  +\frac{e}{c}\mathbf A(\mathbf x)\cdot\dot{\mathbf x}
  -eV(\mathbf x)
  }.
\end{equation}
В лекции 10 будут применены уравнения Эйлера--Лагранжа и явно показано, что это выражение дает силу Лоренца.

\begin{figure}[htbp]
\centering
\includegraphics[width=0.82\textwidth]{lecture_09_figure_09.png}
\caption{Завершенное построение действия на доске: свободный кинетический член плюс \((e/c)A_i\dot x_i\).}
\lectureframe{01:26:36}
\end{figure}

\section{Калибровочные преобразования меняют только конечные точки}

Магнитный вклад в действие равен
\begin{equation}
  S_{\mathrm{mag}}
  =
  \frac{e}{c}
  \int_{t_i}^{t_f}A_i(\mathbf x)\dot x_i\,dt.
\end{equation}
Поскольку \(\dot x_i\,dt=dx_i\), это криволинейный интеграл вдоль пути частицы \(\gamma\):
\begin{equation}
  S_{\mathrm{mag}}
  =
  \frac{e}{c}\int_\gamma A_i\,dx_i
  =
  \frac{e}{c}\int_\gamma\mathbf A\cdot d\mathbf x.
\end{equation}
Разбейте путь на короткие отрезки, возьмите компоненту \(\mathbf A\) вдоль каждого отрезка, умножьте на длину отрезка и сложите. Таков геометрический смысл этого интеграла.

Теперь выполним калибровочное преобразование
\begin{equation}
  A_i\longmapsto A_i+\partial_i s.
\end{equation}
Действие изменяется на
\begin{align}
  \Delta S
  &=
  \frac{e}{c}\int_\gamma \partial_i s\,dx_i\\
  &=
  \frac{e}{c}\int_\gamma ds\\
  &=
  \frac{e}{c}
  \left[
    s(\mathbf x_f)-s(\mathbf x_i)
  \right].
\end{align}
Каждый бесконечно малый отрезок дает значение \(s\) на одном конце минус его значение на другом. При суммировании все внутренние вклады телескопически сокращаются, и остаются только две конечные точки.

\begin{figure}[htbp]
\centering
\includegraphics[width=0.82\textwidth]{lecture_09_figure_10.png}
\caption{Криволинейный интеграл, его разбиение на малые участки пути и результат на концах \(s(\mathbf x_f)-s(\mathbf x_i)\).}
\lectureframe{01:32:31}
\end{figure}

\begin{classroomqa}
\audiencequestion{Если действие меняется при калибровочном преобразовании, почему траектория не меняется?}
\lecturerresponse{Принцип действия сравнивает пути при фиксированных начальной и конечной точках. Добавляемая величина зависит только от этих концов, поэтому для любого пути в данном сравнении она имеет одно и то же значение. Ее вариация равна нулю. Численное значение действия не является строго инвариантным, но уравнения Эйлера--Лагранжа и физическая траектория калибровочно инвариантны.}
\end{classroomqa}

Это различие — главный концептуальный результат лекции:
\begin{equation}
  \begin{aligned}
    S&\ \hbox{may shift by an endpoint term},\\
    \delta S&\ \hbox{and the equations of motion do not}.
  \end{aligned}
\end{equation}
Следовательно, калибровочная свобода совместима с принципом действия. Векторный потенциал может входить в лагранжиан, хотя магнитное поле не выделяет единственного \(\mathbf A\).

\section{Как продолжается рассуждение}

В этой лекции собраны необходимые средства и само действие, но она намеренно останавливается до варьирования этого действия. Следующие шаги теперь определены:
\begin{enumerate}
  \item применить уравнения Эйлера--Лагранжа и получить
  \(\mathbf F_{\mathrm{mag}}=(e/c)\mathbf v\times\mathbf B\);
  \item вычислить канонический импульс, который будет отличаться от \(m\mathbf v\);
  \item построить гамильтониан;
  \item выявить законы сохранения;
  \item решить задачу о движении заряженной частицы в однородном магнитном поле.
\end{enumerate}

Путь к этим результатам объясняет, зачем был введен векторный потенциал. Магнитное поле — это локально измеримый ротор, но векторный потенциал — объект, который естественным образом сопрягается с траекторией частицы. Калибровочные преобразования изменяют это сопряжение только на фиксированных концах, оставляя механику неизменной.
